\documentclass[11pt]{article}
\usepackage[margin=1in]{geometry}
\usepackage{amsmath,amssymb,amsthm}
\usepackage{algorithm}
\usepackage{algpseudocode}
\usepackage{booktabs}
\usepackage{hyperref}

\newtheorem{definition}{Definition}
\newtheorem{theorem}{Theorem}
\newtheorem{lemma}[theorem]{Lemma}
\newtheorem{property}{Property}

\title{\textbf{Mathematical Specification, Taylor Truncation Bounds, and Relative Error Metrics of Finite Difference Gradient Checking (ALGO-NN-26)}}
\author{Team Scaibu \\ \small Email: \texttt{jaydeep.9664920749@gmail.com} \\ \small Architecture and Core Engine Team}
\date{\today}

\begin{document}
\maketitle

\begin{abstract}
Gradient checking provides independent verification of backpropagation implementations by comparing analytic gradients $\mathbf{g}_{\mathrm{analytic}}$ against numerical central finite differences. We establish the quadratic truncation error bound $\mathcal{O}(\epsilon^2)$, define scale-invariant relative error metrics, provide complete pseudocode matching the reference implementation, and derive complexity.
\end{abstract}

\section{Notation and Mathematical Preliminaries}
\label{sec:notation}

Let $\theta \in \mathbb{R}^P$ denote parameter coordinates and $\mathcal{L}: \mathbb{R}^P \to \mathbb{R}$ a scalar loss functional. Let $\mathbf{e}_i \in \mathbb{R}^P$ denote the $i$-th standard basis vector.

\subsection{Central Difference and Relative Error Formulations}
\paragraph{Intuitive Meaning:} Central differences approximate the tangent slope by evaluating the objective symmetrically perturbed by $+\epsilon$ and $-\epsilon$. Comparing this slope against the backpropagated derivative reveals subtle coding errors or broken graph detachments.

\paragraph{Formal Mathematical Definition:}
\begin{equation}
\label{eq:central_diff}
g_{\mathrm{num}, i} = \frac{\mathcal{L}(\theta + \epsilon \mathbf{e}_i) - \mathcal{L}(\theta - \epsilon \mathbf{e}_i)}{2\epsilon}.
\end{equation}
The scale-invariant relative error metric is:
\begin{equation}
\label{eq:rel_error}
\operatorname{RelError}(g_{\mathrm{analytic}, i}, g_{\mathrm{num}, i}) = \frac{|g_{\mathrm{analytic}, i} - g_{\mathrm{num}, i}|}{\max(|g_{\mathrm{analytic}, i}|, |g_{\mathrm{num}, i}|, 10^{-8})}.
\end{equation}

\paragraph{Worked Numerical Example:}
Let analytic gradient $g_{\mathrm{ana}} = 1.000000$, and perturbed losses $\mathcal{L}(\theta + 10^{-6}) = 2.000001$, $\mathcal{L}(\theta - 10^{-6}) = 1.999999$:
$g_{\mathrm{num}} = \frac{2.000001 - 1.999999}{2 \times 10^{-6}} = \frac{0.000002}{0.000002} = 1.000000$.
$\operatorname{RelError} = \frac{|1.0 - 1.0|}{\max(1.0, 1.0, 10^{-8})} = 0.0 \le 10^{-5}$ (Valid).

\subsection{Summary of Symbols}
\begin{itemize}
    \item $P \in \mathbb{N}$: Number of checked parameters.
    \item $\boldsymbol{\theta} \in \mathbb{R}^P$: Parameter vector.
    \item $\mathbf{g}_{\mathrm{ana}} \in \mathbb{R}^P$: Backpropagated analytic gradient vector.
    \item $\mathbf{g}_{\mathrm{num}} \in \mathbb{R}^P$: Numerical finite difference gradient vector.
    \item $\epsilon \in \mathbb{R}_{>0}$: Perturbation step ($10^{-6}$).
    \item $\tau_{\mathrm{tol}} \in \mathbb{R}_{>0}$: Validation relative error threshold ($10^{-5}$).
\end{itemize}

\section{Problem Definition and Core Concepts}
\label{sec:problem_definition}

\subsection{Taylor Expansion Error Cancellation}
Expanding $\mathcal{L}(\theta \pm \epsilon)$ via Taylor's theorem:
\begin{equation}
\mathcal{L}(\theta + \epsilon) = \mathcal{L}(\theta) + \epsilon \mathcal{L}'(\theta) + \frac{\epsilon^2}{2} \mathcal{L}''(\theta) + \frac{\epsilon^3}{6} \mathcal{L}'''(\theta) + \mathcal{O}(\epsilon^4).
\end{equation}
\begin{equation}
\mathcal{L}(\theta - \epsilon) = \mathcal{L}(\theta) - \epsilon \mathcal{L}'(\theta) + \frac{\epsilon^2}{2} \mathcal{L}''(\theta) - \frac{\epsilon^3}{6} \mathcal{L}'''(\theta) + \mathcal{O}(\epsilon^4).
\end{equation}
Subtracting eliminates all even-order terms, proving central differences achieve $\mathcal{O}(\epsilon^2)$ truncation accuracy compared to forward differences $\mathcal{O}(\epsilon)$.

\section{Algorithmic Specification}
\label{sec:algorithm}

Algorithm~\ref{alg:grad_check} specifies the evaluation routine matching \texttt{impl.py}.

\begin{algorithm}[H]
\caption{Finite Difference Gradient Checking Routine}
\label{alg:grad_check}
\begin{algorithmic}[1]
\Require Parameters $\boldsymbol{\theta} \in \mathbb{R}^P$, analytic gradients $\mathbf{g}_{\mathrm{ana}} \in \mathbb{R}^P$, perturbed losses $\mathbf{L}_+, \mathbf{L}_- \in \mathbb{R}^P$, step $\epsilon > 0$, tolerance $\tau_{\mathrm{tol}} > 0$.
\Ensure Status $\mathrm{is\_correct}$, maximum error $E_{\max}$, numerical gradients $\mathbf{g}_{\mathrm{num}}$, relative errors $\mathbf{e}_{\mathrm{rel}}$.
\State Initialize $\mathbf{g}_{\mathrm{num}} \gets [], \mathbf{e}_{\mathrm{rel}} \gets [], E_{\max} \gets 0.0, \mathrm{passed} \gets \text{True}$
\For{$i = 1$ to $P$}
    \State $g_n \gets \frac{L_{+, i} - L_{-, i}}{2.0 \cdot \epsilon}$
    \State $g_a \gets g_{\mathrm{ana}, i}$
    \State $diff \gets |g_a - g_n|$
    \State $scale \gets \max(|g_a|, |g_n|, 10^{-8})$
    \State $err \gets \frac{diff}{scale}$
    \State Append $g_n$ to $\mathbf{g}_{\mathrm{num}}$, append $err$ to $\mathbf{e}_{\mathrm{rel}}$
    \If{$err > E_{\max}$}
        \State $E_{\max} \gets err$
    \EndIf
    \If{$err > \tau_{\mathrm{tol}}$}
        \State $\mathrm{passed} \gets \text{False}$
    \EndIf
\EndFor
\State \Return $\mathrm{passed}, \quad E_{\max}, \quad \mathbf{g}_{\mathrm{num}}, \quad \mathbf{e}_{\mathrm{rel}}$
\end{algorithmic}
\end{algorithm}

\section{Theoretical Guarantees and Error Bounds}
\label{sec:guarantees}

\begin{theorem}[Quadratic Truncation Error Bound]
\label{thm:taylor_quadratic_error}
\textbf{Assumptions:} Let $\mathcal{L}$ be three times continuously differentiable ($C^3$) with $|\mathcal{L}'''(\theta)| \le M_3$.
\textbf{Guarantee:} The central finite difference approximation satisfies $|g_{\mathrm{num}} - g_{\mathrm{analytic}}| \le \frac{M_3}{6} \epsilon^2$.
\textbf{Proof:}
From the Taylor remainder theorem:
$\mathcal{L}(\theta + \epsilon) - \mathcal{L}(\theta - \epsilon) = 2\epsilon \mathcal{L}'(\theta) + \frac{\epsilon^3}{3} \mathcal{L}'''(\xi)$ for some $\xi \in [\theta - \epsilon, \theta + \epsilon]$.
Dividing by $2\epsilon$:
$\frac{\mathcal{L}(\theta + \epsilon) - \mathcal{L}(\theta - \epsilon)}{2\epsilon} - \mathcal{L}'(\theta) = \frac{\epsilon^2}{6} \mathcal{L}'''(\xi) \implies |g_{\mathrm{num}} - g_{\mathrm{analytic}}| \le \frac{M_3}{6} \epsilon^2$.
\textbf{Limitations:} If $\epsilon$ is chosen too small ($\epsilon < 10^{-12}$ in FP64), floating-point subtractive cancellation errors dominate.
\end{theorem}

\section{Complexity Analysis}
\label{sec:complexity}

\begin{itemize}
    \item \textbf{Time Complexity:} $\mathcal{O}(P)$ operations for finite difference estimation given perturbed values. Evaluating perturbed values requires $2P$ forward passes.
    \item \textbf{Space Complexity:} $\mathcal{O}(P)$ storage for numerical gradients and error vectors.
\end{itemize}

\section{Worked Numerical Example}
\label{sec:worked_example}

Let parameter $\theta = [0.5]$, analytic gradient $g_{\mathrm{ana}} = 2.0$, perturbed losses $L_+ = 1.000002000001, L_- = 0.999998000001$ with $\epsilon = 10^{-6}$:
\begin{enumerate}
    \item $L_+ - L_- = 0.000004000000$.
    \item $g_{\mathrm{num}} = \frac{0.000004000000}{2 \times 10^{-6}} = 2.000000$.
    \item Relative error: $\frac{|2.0 - 2.0|}{\max(2.0, 2.0, 10^{-8})} = 0.0 \implies \text{Valid}$.
\end{enumerate}

\section{Numerical Considerations and Edge Cases}
\label{sec:numerical}

\begin{itemize}
    \item \textbf{Non-Smooth Kinks (ReLU at 0):} At non-differentiable points (such as $\operatorname{ReLU}(0)$), analytic subgradients and finite difference secants may disagree; test points should be sampled away from kinks.
\end{itemize}

\begin{thebibliography}{9}
\bibitem{nocedal2006} J.~Nocedal and S.~J.~Wright, \emph{Numerical Optimization}, 2nd~ed., Springer, 2006.
\end{thebibliography}

\end{document}
